% 分章页86—90附录2.A末；印刷页158—162。
平滑样条回归函数为
\begin{equation}
\widehat f(\boldsymbol x)=\sum_{j=1}^r\widehat a_j\phi_j(\boldsymbol x)
 +\sum_{i=1}^n\widehat b_i\xi_i(\boldsymbol x),\tag{2.A.105}
\end{equation}
其中 $\xi_i(\boldsymbol x)=L_{i(z)}R_1(x,z)$。最后，得到拟合值
\begin{equation}
\widehat f=H(\lambda)Y,\tag{2.A.106}
\end{equation}
其中 $H(\lambda)=I-n\lambda Q_2(Q_2^TMQ_2)^{-1}Q_2^T$。

\begin{bookexample}{2.A.5：$W_2^m[a,b]\otimes W_2^m[a,b]$ 上的平滑样条}
考虑区间 $[0,1]$ 上的两个连续变量 $x_1$ 与 $x_2$。假设这两个变量的模型空间为张量积 $W_2^m[a,b]\otimes W_2^m[a,b]$。设 $m_1=m_2=2$，即考虑三次样条。

定义四个平均算子：
\[
A_1^{(k)}f=\int_0^1f\,dx_k,
\qquad
A_2^{(k)}f=\int_0^1\frac{\partial f(\boldsymbol x)}{\partial x_k}
 \,dx_k(x_k-0.5),\qquad k=1,2.
\]
因此，有
\begingroup
\allowdisplaybreaks[2]
\begin{align}
A_1^{(1)}A_1^{(2)}f
 &=A_1^{(1)}\left(\int_0^1f(x_1,x_2)\,dx_2\right)\notag\\
 &=\int_0^1\int_0^1f(x_1,x_2)\,dx_1dx_2=\mu,\notag\\
A_2^{(1)}A_1^{(2)}f
 &=A_2^{(1)}\left(\int_0^1f(x_1,x_2)\,dx_2\right)\notag\\
 &=\int_0^1\int_0^1\frac{\partial f(x_1,x_2)}{\partial x_1}
    \,dx_1dx_2(x_1-0.5)\notag\\
 &=\beta_1(x_1-0.5),\notag\\
A_1^{(1)}A_2^{(2)}f
 &=A_1^{(1)}\left(\int_0^1
   \frac{\partial f(\boldsymbol x)}{\partial x_2}
   \,dx_2(x_2-0.5)\right)\notag\\
 &=\int_0^1\int_0^1\frac{\partial f(\boldsymbol x)}{\partial x_2}
   \,dx_1dx_2(x_2-0.5)\notag\\
 &=\beta_2(x_2-0.5),\notag\\
A_2^{(1)}A_2^{(2)}f
 &=A_2^{(1)}\left(\int_0^1
   \frac{\partial f(\boldsymbol x)}{\partial x_2}
   \,dx_2(x_2-0.5)\right)\notag\\
 &=\int_0^1\int_0^1
   \frac{\partial^2 f(\boldsymbol x)}{\partial x_1\,\partial x_2}
   \,dx_1dx_2(x_1-0.5)(x_2-0.5)\notag\\
 &=\gamma_{12}(x_1-0.5)(x_2-0.5),\tag{2.A.107}
\end{align}
\endgroup
其中
\[
\begin{aligned}
\mu&=\int_0^1\int_0^1f(x_1,x_2)\,dx_1dx_2,\\
\beta_1&=\int_0^1\int_0^1
       \frac{\partial f(x_1,x_2)}{\partial x_1}\,dx_1dx_2,\\
\beta_2&=\int_0^1\int_0^1
       \frac{\partial f(\boldsymbol x)}{\partial x_2}\,dx_1dx_2,\\
\gamma_{12}&=\int_0^1\int_0^1
       \frac{\partial^2f(\boldsymbol x)}{\partial x_1\,\partial x_2}
       \,dx_1dx_2.
\end{aligned}
\]
定义 $A_3^{(k)}=I-A_1^{(k)}-A_2^{(k)}$，$k=1,2$，则有
\[
A_3^{(k)}f=f(\boldsymbol x)-\int_0^1f\,dx_k
 -\int_0^1\frac{\partial f(\boldsymbol x)}{\partial x_k}
     \,dx_k(x_k-0.5),
\]
\[
\begin{aligned}
A_1^{(1)}A_3^{(2)}f
 &=\int_0^1f(x_1,x_2)\,dx_1
   -\int_0^1\int_0^1f(x_1,x_2)\,dx_1dx_2\\
 &\quad-\int_0^1\int_0^1
    \frac{\partial f(x_1,x_2)}{\partial x_2}\,dx_1dx_2(x_2-0.5)\\
 &=g_2(x_2)-\mu-\beta_2(x_2-0.5),
\end{aligned}
\]
\[
\begin{aligned}
A_2^{(1)}A_3^{(2)}f
 &=\int_0^1\frac{\partial f(x_1,x_2)}{\partial x_1}\,dx_1(x_1-0.5)\\
 &\quad-\int_0^1\int_0^1
    \frac{\partial f(x_1,x_2)}{\partial x_1}\,dx_1dx_2(x_1-0.5)\\
 &\quad-\int_0^1\int_0^1
    \frac{\partial^2f(x_1,x_2)}{\partial x_1\,\partial x_2}
    \,dx_1dx_2(x_1-0.5)(x_2-0.5)\\
 &=\alpha_1(x_1-0.5)-\beta_1(x_1-0.5)
   -\gamma_{12}(x_1-0.5)(x_2-0.5),
\end{aligned}
\]
\[
\begin{aligned}
A_3^{(1)}A_1^{(2)}f
 &=\int_0^1f(x_1,x_2)\,dx_2
  -\int_0^1\int_0^1f(x_1,x_2)\,dx_1dx_2\\
 &\quad-\int_0^1\int_0^1
    \frac{\partial f(x_1,x_2)}{\partial x_1}\,dx_1dx_2(x_1-0.5)\\
 &=g_1(x_1)-\mu-\beta_1(x_1-0.5),
\end{aligned}
\]
\[
\begin{aligned}
A_3^{(1)}A_2^{(2)}f
 &=\int_0^1\frac{\partial f(x_1,x_2)}{\partial x_2}\,dx_2(x_2-0.5)\\
 &\quad-\int_0^1\int_0^1
    \frac{\partial f(x_1,x_2)}{\partial x_2}\,dx_1dx_2(x_2-0.5)\\
 &\quad-\int_0^1\int_0^1
    \frac{\partial^2f(x_1,x_2)}{\partial x_1\,\partial x_2}
    \,dx_1dx_2(x_1-0.5)(x_2-0.5)\\
 &=\alpha_2(x_2-0.5)-\beta_2(x_2-0.5)
   -\gamma_{12}(x_1-0.5)(x_2-0.5),
\end{aligned}
\]
\[
\begin{aligned}
A_3^{(1)}A_3^{(2)}f
 &=f(x_1,x_2)+\int_0^1\int_0^1f(x_1,x_2)\,dx_1dx_2\\
 &\quad-\int_0^1f(x_1,x_2)\,dx_1-\int_0^1f(x_1,x_2)\,dx_2\\
 &\quad+\left[\int_0^1\int_0^1
    \frac{\partial f(x_1,x_2)}{\partial x_1}\,dx_1dx_2
    -\int_0^1\frac{\partial f(x_1,x_2)}{\partial x_1}\right](x_1-0.5)\\
 &\quad+\left[\int_0^1\int_0^1
    \frac{\partial f(x_1,x_2)}{\partial x_2}\,dx_1dx_2\right.\\
 &\quad-\int_0^1\frac{\partial f(x_1,x_2)}{\partial x_2}
    \,dx_2(x_2-0.5)\\
 &\quad+\int_0^1\int_0^1
    \frac{\partial^2f(x_1,x_2)}{\partial x_1\,\partial x_2}
    \,dx_1dx_2(x_1-0.5)(x_2-0.5)\\
 &=f(x_1,x_2)+\mu-g_1(x_1)-g_2(x_2)
   +(\beta_1-\alpha_1)(x_1-0.5)\\
 &\quad+(\beta_2-\alpha_2)(x_2-0.5)
   +\gamma_{12}(x_1-0.5)(x_2-0.5),
\end{aligned}
\]
其中
\[
\begin{aligned}
g_1(x_1)&=\int_0^1f(x_1,x_2)\,dx_2,&
g_2(x_2)&=\int_0^1f(x_1,x_2)\,dx_1,\\
\alpha_1&=\int_0^1\frac{\partial f(x_1,x_2)}{\partial x_1}\,dx_1,&
\alpha_2&=\int_0^1\frac{\partial f(x_1,x_2)}{\partial x_2}\,dx_2.
\end{aligned}
\]
利用式（2.A.80），函数 $f(x_1,x_2)$ 在张量积空间 $W_2^m[a,b]\otimes W_2^m[a,b]$ 上的分解为
\[
\begin{aligned}
f&=\left\{A_1^{(1)}+A_2^{(1)}+A_3^{(1)}\right\}
   \left\{A_1^{(2)}+A_2^{(2)}+A_3^{(2)}\right\}f\\
 &=A_1^{(1)}A_1^{(2)}f+A_1^{(1)}A_2^{(2)}f+A_1^{(1)}A_3^{(2)}f\\
 &\quad+A_2^{(1)}A_1^{(2)}f+A_2^{(1)}A_2^{(2)}f+A_2^{(1)}A_3^{(2)}f\\
 &\quad+A_3^{(1)}A_1^{(2)}f+A_3^{(1)}A_2^{(2)}f+A_3^{(1)}A_3^{(2)}f\\
 &=\mu+\beta_2(x_2-0.5)+f_2^s(x_2)+\beta_1(x_1-0.5)\\
 &\quad+\gamma_{12}(x_1-0.5)(x_2-0.5)
   +f_{12}^{ls}(x_1,x_2)+f_1^s(x_1)\\
 &\quad+f_{12}^{sl}(x_1,x_2)+f_{12}^{ss}(x_1,x_2),
\end{aligned}
\]
其中
\[
f_1^s(x_1)=A_3^{(1)}A_1^{(2)}f
 =g_1(x_1)-\mu-\beta_1(x_1-0.5),
\]
\[
f_2^s(x_2)=A_1^{(1)}A_3^{(2)}f
 =g_2(x_2)-\mu-\beta_2(x_2-0.5),
\]
\[
\begin{aligned}
f_{12}^{ls}(x_1,x_2)&=A_2^{(1)}A_3^{(2)}f\\
 &=\alpha_1(x_1-0.5)-\beta_1(x_1-0.5)
   -\gamma_{12}(x_1-0.5)(x_2-0.5),
\end{aligned}
\]
\[
\begin{aligned}
f_{12}^{sl}(x_1,x_2)&=A_3^{(1)}A_2^{(2)}f\\
 &=\alpha_2(x_2-0.5)-\beta_2(x_2-0.5)
   -\gamma_{12}(x_1-0.5)(x_2-0.5),
\end{aligned}
\]
\[
\begin{aligned}
f_{12}^{ss}(x_1,x_2)&=A_3^{(1)}A_3^{(2)}f\\
 &=f(x_1,x_2)+\mu-g_1(x_1)-g_2(x_2)
   +(\beta_1-\alpha_1)(x_1-0.5)\\
 &\quad+(\beta_2-\alpha_2)(x_2-0.5)
   +\gamma_{12}(x_1-0.5)(x_2-0.5).
\end{aligned}
\]
$f_1^s(x_1)$ 与 $f_2^s(x_2)$ 分别度量 $x_1$ 与 $x_2$ 的主效应；$f_{12}^{ls}(x_1,x_2)$、$f_{12}^{sl}(x_1,x_2)$ 和 $f_{12}^{ss}(x_1,x_2)$ 分别为 $x_1$ 与 $x_2$ 的线性—平滑、平滑—线性和平滑—平滑交互作用函数。式（2.A.27）说明，函数可以分解为线性函数和平滑的非线性交互作用。

接着研究模型空间 $W_2^m[0,1]\otimes W_2^m[0,1]$ 的张量积分解。假设 Sobolev 空间 $W_2^m[0,1]$ 分解为由两个基函数构成的有限维基函数子空间 $H_0^{(k)}=\{1\}$、$H_1^{(k)}=\{x_k-0.5\}$，以及无限维函数空间
\[
H_2^{(k)}=\left\{f\in W_2^m[0,1]:
 \int_0^1f(x_1,x_2)\,dx_k=0,\
 \int_0^1\frac{\partial f(\boldsymbol x)}{\partial x_k}\,dx_k=0\right\},
\quad k=1,2.
\]
张量积空间 $W_2^m[0,1]\otimes W_2^m[0,1]$ 可分解为
\begin{equation}
\begin{aligned}
&W_2^m[0,1]\otimes W_2^m[0,1]\\
&=\left\{H_0^{(1)}\oplus H_1^{(1)}\oplus H_2^{(1)}\right\}
 \left\{H_0^{(2)}\oplus H_1^{(2)}\oplus H_2^{(2)}\right\}\\
&=\left\{H_0^{(1)}\otimes H_0^{(2)}\right\}
 \oplus\left\{H_0^{(1)}\otimes H_1^{(2)}\right\}
 \oplus\left\{H_0^{(1)}\otimes H_2^{(2)}\right\}\\
&\quad\oplus\left\{H_1^{(1)}\otimes H_0^{(2)}\right\}
 \oplus\left\{H_1^{(1)}\otimes H_1^{(2)}\right\}
 \oplus\left\{H_1^{(1)}\otimes H_2^{(2)}\right\}\\
&\quad\oplus\left\{H_2^{(1)}\otimes H_0^{(2)}\right\}
 \oplus\left\{H_2^{(1)}\otimes H_1^{(2)}\right\}
 \oplus\left\{H_2^{(1)}\otimes H_2^{(2)}\right\}.
\end{aligned}\tag{2.A.108}
\end{equation}
式（2.A.28）可以重新分组为
\begin{equation}
\begin{aligned}
W_2^m[0,1]\otimes W_2^m[0,1]
 &=H^{(0)}\oplus H^{(1)}\oplus H^{(2)}\\
 &\quad\oplus H^{(3)}\oplus H^{(4)}\oplus H^{(5)},
\end{aligned}\tag{2.A.109}
\end{equation}
其中
\[
\begin{aligned}
H^{(0)}
 &=\left\{H_0^{(1)}\otimes H_0^{(2)}\right\}
 \oplus\left\{H_1^{(1)}\otimes H_0^{(2)}\right\}\\
 &\quad\oplus\left\{H_0^{(1)}\otimes H_1^{(2)}\right\}
 \oplus\left\{H_1^{(1)}\otimes H_1^{(2)}\right\}\\
 &=\{\phi_1(\boldsymbol x)=1\}
 \oplus\{\phi_2(\boldsymbol x)=x_1-0.5\}\\
 &\quad\oplus\{\phi_3(\boldsymbol x)=x_2-0.5\}
 \oplus\{\phi_4(\boldsymbol x)=(x_1-0.5)(x_2-0.5)\},
\end{aligned}
\]
\[
\begin{aligned}
H^1&=H_2^{(1)}\otimes H_0^{(2)},&
H^2&=H_2^{(1)}\otimes H_1^{(2)},\\
H^3&=H_0^{(1)}\otimes H_2^2,&
H^4&=H_1^{(1)}\otimes H_2^{(2)},\\
H^5&=H_2^{(1)}\otimes H_2^{(2)}.
\end{aligned}
\]
$H^1,H^2,H^3,H^4,H^5$ 的再生核 $R^1,R^2,R^3,R^4,R^5$ 分别计算如下：
\begin{equation}
\begin{aligned}
R^1&=k_2(x_1)k_2(z_1)-k_4(|x_1-z_1|),\\
R^2&=\bigl(k_2(x_1)k_2(z_1)-k_4(|x_1-z_1|)\bigr)k_1(x_2)k_1(z_2),\\
R^3&=k_2(x_2)k_2(z_2)-k_4(|x_2-z_2|),\\
R^4&=k_1(x_1)k_1(z_1)\bigl(k_2(x_2)k_2(z_2)-k_4(|x_2-z_2|)\bigr),\\
R^5&=\bigl(k_2(x_1)k_2(z_1)-k_4(|x_1-z_1|)\bigr)\\
 &\quad{}\times\bigl(k_2(x_2)k_2(z_2)-k_4(|x_2-z_2|)\bigr).
\end{aligned}\tag{2.A.110}
\end{equation}
假设连续泛函为求值泛函，则矩阵 $T$ 为
\begin{equation}
T=\begin{bmatrix}
1&x_1^{(1)}-0.5&x_2^{(1)}-0.5&(x_1^{(1)}-0.5)(x_2^{(1)}-0.5)\\
1&x_1^{(2)}-0.5&x_2^{(2)}-0.5&(x_1^{(2)}-0.5)(x_2^{(2)}-0.5)\\
\vdots&\vdots&\vdots&\vdots\\
1&x_1^{(n)}-0.5&x_2^{(n)}-0.5&(x_1^{(n)}-0.5)(x_2^{(n)}-0.5)
\end{bmatrix},\tag{2.A.111}
\end{equation}
其中 $\boldsymbol x_1=[x_1^{(1)},\ldots,x_1^{(n)}]^T$ 和 $\boldsymbol x_2=[x_2^{(1)},\ldots,x_2^{(n)}]^T$ 是两个观测自变量向量。

接着计算矩阵 $\Sigma_j$：
\[
\Sigma_1=\begin{bmatrix}
R^1(\boldsymbol x^{(1)},\boldsymbol x^{(1)})&\cdots&
 R^1(\boldsymbol x^{(1)},\boldsymbol x^{(n)})\\
\vdots&\vdots&\vdots\\
R^1(\boldsymbol x^{(n)},\boldsymbol x^{(1)})&\cdots&
 R^1(\boldsymbol x^{(n)},\boldsymbol x^{(n)})
\end{bmatrix}
\]
\[
=\begin{bmatrix}
[k_2(x_1^{(1)})]^2-k_4(0)&\cdots&
 k_2(x_1^{(1)})k_2(x_1^{(n)})-k_4(|x_1^{(1)}-x_1^{(n)}|)\\
\vdots&\vdots&\vdots\\
k_2(x_1^{(n)})k_2(x_1^{(1)})-k_4(|x_1^{(n)}-x_1^{(1)}|)&\cdots&
 [k_2(x_1^{(n)})]^2k_4(0)
\end{bmatrix}.
\]
类似地，利用式（2.A.12）可以计算矩阵 $\Sigma_2,\Sigma_3,\Sigma_4,\Sigma_5$。最后，计算
\[
\Sigma_\theta=\theta_1\Sigma_1+\theta_2\Sigma_2+\theta_3\Sigma_3
 +\theta_4\Sigma_4+\theta_5\Sigma_5.
\]
\end{bookexample}
